\mysection{Background}
\label{sec:background}

\MyPara{SLOs of Online LLM Serving.}
In online LLM serving, unlike offline batch processing, the system must interact with users or downstream applications in real-time. 
Consequently, providers enforce strict SLOs to guarantee the quality of experience. 
Two key latency metrics define these objectives. 
First, \textit{Time-To-First-Token} (TTFT) measures the latency to process the input prompt and generate the first token, which is a crucial metric for interactive applications, such as chatbots. 
Second, \textit{Time-Per-Output-Token} (TPOT) measures the average inter-token latency during the autoregressive generation, which is essential to ensure that text generation is fluid and matches human reading speeds.

\MyPara{Memory footprint of LLM serving.}
LLM inference is inherently memory-bound, with GPU memory dominated by model weights and KV cache.

\noindent \textbf{\textit{Model weights.}}
LLM weights are massive; for example, a 70B-parameter model requires roughly 140~GB of GPU memory in FP16. 
Under long-tail workloads~\cite{aegaeon}, idle models waste GPU memory; prior work~\cite{aegaeon,fu2024serverlessllm} mitigates this with model eviction and accelerated weight loading.

\noindent \textbf{\textit{KV cache and PagedAttention}.}
The KV cache stores intermediate Key (K) and Value (V) tensors for attention layers to avoid redundant computation. 
Unlike weights, the KV cache is dynamic, growing with sequence length and batch size. 

To improve GPU memory efficiency, mainstream LLM serving systems adopt PagedAttention~\cite{kwon2023efficient}, which splits each request’s KV cache into fixed-size blocks that can be stored non-contiguously and reused across requests.
Frameworks such as SGLang~\cite{sglang@arxiv23} and vLLM~\cite{kwon2023efficient} implement this design by reserving large PyTorch tensors as a KV block pool and mapping each request’s logical blocks to physical offsets in the pool. For each of the $L$ attention layers, a K tensor and a V tensor of shape $(B, T, H, D)$ are allocated, where $B$ is the number of blocks, $T$ is the tokens per block, $H$ is the number of attention heads, and $D$ is the head dimension. The values of $L$, $H$, and $D$ vary across model architectures (\eg, Llama-3-8B ($32, 8, 128$) vs. GPT-OSS-120B ($36, 8, 64$)), resulting in different KV block sizes.

